\section{An exact torus twist calculation and a constant orbit series}
\label{sec:exacttwist}

\subsection{The length function and its integral}

Let $X$ be a one-holed torus with boundary length $b\ge0$. Cut along a
simple curve $\alpha$ of length $a>0$, and choose a curve $\beta$ crossing
$\alpha$ once. Normalize the twist origin so that this crossing curve
has minimum length at $\tau=0$. Set
\begin{equation}
 A^2=\frac{\cosh a+\cosh(b/2)}{\cosh a-1}.
 \label{eq:Atwist}
\end{equation}
The exact formula is
\begin{equation}
 \boxed{\cosh\bigl(\ell_\beta(\tau)/2\bigr)=A\cosh(\tau/2).}
 \label{eq:torustwist}
\end{equation}
It follows from the right-angled-hexagon description of the cut pants;
see also~\cite[Proposition 3.1]{ALPS}. For $b=0$,
$A=\coth(a/2)$. The minimum length is $r_0=2\arcosh A$.

For any nonnegative measurable function $f$, change variables separately
on the positive and negative twist axes to obtain
\begin{equation}
 \int_{\RR}f\bigl(\ell_\beta(u)\bigr)\,\dd u
 =2\int_{r_0}^\infty
 f(r)\frac{\sinh(r/2)}{\sqrt{\cosh^2(r/2)-A^2}}\,\dd r.
 \label{eq:toruscoarea}
\end{equation}
Indeed, differentiating $\cosh(r/2)=A\cosh(u/2)$ gives
$\dd u/\dd r=\sinh(r/2)/\sqrt{\cosh^2(r/2)-A^2}$ for $u>0$.
The two signs of $u$ explain the factor $2$.

Several useful checks are immediate:
\begin{align}
 \int_{\RR}\sech^2\bigl(\ell_\beta(u)/2\bigr)\,\dd u
     &=\frac4{A^2},\label{eq:sechintegral}\\
 \operatorname{meas}\{u:\ell_\beta(u)\le R\}
     &=4\arcosh\left(\frac{\cosh(R/2)}{A}\right),\quad R\ge r_0,
       \label{eq:sublevel}\\
 \int_{\RR}\cosh^{-s}\bigl(\ell_\beta(u)/2\bigr)\,\dd u
     &=\frac{2\sqrt\pi\,\Gamma(s/2)}
             {A^s\Gamma((s+1)/2)},\quad s>0.
       \label{eq:betaintegral}
\end{align}
The sublevel measure is zero for $R<r_0$. These are exact twist integrals,
not asymptotic estimates.

\subsection{A positive partition of unity}

Define
\begin{equation}
 H_a(u)=\frac{\sinh(a/2)}{\cosh u+\cosh(a/2)}
 =\frac12\left[
 \tanh\left(\frac{u+a/2}{2}\right)
 -\tanh\left(\frac{u-a/2}{2}\right)\right].
 \label{eq:Hkernel}
\end{equation}
The second expression telescopes under translation by $a$, so
\begin{equation}
 \sum_{n\in\ZZ}H_a(\tau+na)=1.
 \label{eq:Hidentity}
\end{equation}
All terms are positive and the convergence is exponential on compact
twist intervals.

Since~\eqref{eq:torustwist} implies
$\cosh r+1=A^2(\cosh u+1)$, eliminate $u$ from~\eqref{eq:Hkernel} and set
\begin{equation}
 K_{a,b}(r)=
 \frac{A^2\sinh(a/2)}
      {\cosh r+1+A^2(\cosh(a/2)-1)}.
 \label{eq:Kkernel}
\end{equation}
For $r_n=\ell_{D_\alpha^n\beta}(X)$ we obtain the pointwise identity
\begin{equation}
 \boxed{\sum_{n\in\ZZ}K_{a,b}(r_n)=1.}
 \label{eq:Kidentity}
\end{equation}
This is a constant series proved by telescoping, rather than inferred from
its average. Its kernel depends on the distinguished cuff length $a$ and
the boundary length $b$, and its index set is one twist orbit. It is not
an unrestricted universal sum over every curve on the surface.

There is a simple source of many such examples. If $\rho$ is an even
probability density on $\RR$, then
\[
 H_{a,\rho}(u)=\int_{u-a/2}^{u+a/2}\rho(v)\,\dd v
 \quad\text{satisfies}\quad
 \sum_{n\in\ZZ}H_{a,\rho}(u+na)=1.
\]
The translated intervals partition the line. Evenness makes
$H_{a,\rho}$ a function of $|u|$, hence of the length $r$ through
\eqref{eq:torustwist}. The choice
$\rho(v)=\frac14\sech^2(v/2)$ gives~\eqref{eq:Hkernel}.
This realizes the proposed ``zoo'' of identities in a concrete way, with
the source of constancy and the parameter dependence both visible.

